\chapter{Boas práticas, erros comuns e diretrizes internacionais para SDM}
%================================================

\markboth{Boas práticas em SDM}{Boas práticas em SDM}

\section{Introdução}

A Modelagem de Distribuição de Espécies (Species Distribution Modeling -- SDM)
evoluiu rapidamente nas últimas duas décadas, tornando-se uma das principais
ferramentas utilizadas em ecologia, biogeografia, conservação da biodiversidade,
planejamento ambiental e avaliação dos impactos das mudanças climáticas.

Paralelamente ao desenvolvimento de novos algoritmos e ao crescimento da
disponibilidade de dados ambientais e registros de ocorrência, surgiu uma
demanda crescente por maior rigor metodológico, transparência analítica e
reprodutibilidade científica. Atualmente, estudos de SDM não são avaliados
apenas pela sofisticação dos modelos utilizados, mas também pela qualidade dos
dados, clareza metodológica, robustez das inferências e disponibilidade dos
códigos e dados utilizados.

Diversos trabalhos demonstram que muitos dos erros observados em SDM não estão
relacionados aos algoritmos em si, mas às decisões tomadas durante a seleção de
variáveis ambientais, definição da área acessível ($M$), tratamento do viés
amostral, avaliação da extrapolação e interpretação inadequada das métricas de
desempenho \citep{guisan2000predictive, elith2006novel,
	barve2011crucial, guisan2017habitat, zurell2020standard}.

Ao mesmo tempo, periódicos científicos de alto impacto e organizações
internacionais vêm enfatizando a importância da ciência aberta, da
reprodutibilidade computacional e da adoção dos princípios FAIR para garantir
que os resultados produzidos possam ser verificados, reutilizados e integrados
em futuras sínteses ecológicas.

Esta unidade apresenta um conjunto de recomendações práticas destinadas a
auxiliar pesquisadores e estudantes na construção de modelos mais robustos,
transparentes e ecologicamente consistentes. O objetivo não é apenas evitar
erros comuns, mas promover uma visão crítica da modelagem ecológica como um
processo científico integrado que combina teoria ecológica, conhecimento
naturalístico, estatística, ciência de dados e tomada de decisão para a
conservação da biodiversidade.

\section{Por que modelos falham?}

Uma percepção equivocada frequentemente observada entre iniciantes em SDM é a
ideia de que modelos falham principalmente devido à escolha inadequada do
algoritmo. Embora diferentes algoritmos possuam características distintas,
estudos comparativos demonstram que fatores relacionados à qualidade dos dados,
à definição do problema ecológico e às decisões metodológicas exercem impacto
igual ou superior ao da escolha do algoritmo propriamente dito
\citep{elith2006novel, araujo2007ensemble, guisan2017habitat}.

Modelos podem apresentar elevados valores de AUC e TSS e, ainda assim,
produzir interpretações ecologicamente incorretas. Da mesma forma, modelos
estatisticamente simples podem fornecer inferências robustas quando baseados em
dados adequados, hipóteses bem definidas e procedimentos transparentes.

Nas seções seguintes são apresentados alguns dos erros mais frequentes
identificados na literatura especializada.

%------------------------------------------------
\section{Erro 1 - Utilizar todas as variáveis do WorldClim sem seleção prévia}

Um erro muito comum consiste em utilizar automaticamente todas as variáveis
bioclimáticas disponíveis, especialmente as 19 variáveis do WorldClim, sem
avaliar redundância estatística ou relevância ecológica. Embora essas variáveis
representem diferentes componentes do clima, muitas delas são fortemente
correlacionadas entre si.

A inclusão indiscriminada de variáveis ambientais pode aumentar a
multicolinearidade, dificultar a interpretação ecológica dos coeficientes,
favorecer o sobreajuste e reduzir a capacidade de transferência espacial e
temporal dos modelos \citep{dormann2013collinearity, guisan2017habitat}.

\begin{sdmwarningbox}
	Utilizar muitas variáveis ambientais sem seleção prévia pode gerar
	multicolinearidade, sobreajuste, instabilidade nas projeções futuras,
	baixa transferibilidade e interpretações ecológicas inconsistentes.
\end{sdmwarningbox}

Boas práticas incluem a avaliação da correlação entre variáveis, o cálculo do
Fator de Inflação da Variância (VIF), o uso criterioso de PCA quando
apropriado e, principalmente, a seleção de preditores com base em hipóteses
ecológicas relacionadas à espécie e à escala do estudo.

%------------------------------------------------
\section{Erro 2 - Ignorar a área acessível \texorpdfstring{($M$)}{M}}

A definição da área acessível ($M$) é uma das decisões mais importantes em
SDM. No contexto do \textit{BAM framework}, a distribuição observada de uma
espécie resulta da interação entre condições abióticas adequadas, interações
bióticas e áreas historicamente acessíveis por dispersão
\citep{soberon2005interpretation, soberon2009niches, barve2011crucial}.

Ignorar a área $M$ pode levar à seleção de background ou pseudo-ausências em
regiões que a espécie nunca poderia ter alcançado, confundindo ausência de
dispersão com inadequação ambiental.

\begin{sdmwarningbox}
	\textbf{Consequências principais.}
	
	A ausência de uma área $M$ ecologicamente justificada pode gerar
	pseudo-ausências irreais, background inadequado, viés na calibração do modelo
	e superestimação da distribuição potencial.
\end{sdmwarningbox}

A área $M$ deve ser definida com base no conhecimento biogeográfico da espécie,
barreiras geográficas, história de dispersão, limites de biomas, bacias
hidrográficas, ecorregiões ou outras unidades espacialmente coerentes.

%------------------------------------------------
\section{Erro 3 - Escolher o melhor modelo apenas pelo AUC}

O AUC é uma das métricas mais utilizadas em SDM por medir a capacidade
discriminatória do modelo. No entanto, seu uso isolado pode ser problemático,
pois altos valores de AUC não garantem que o modelo seja ecologicamente
plausível ou adequado para projeções espaciais e temporais
\citep{lobo2008auc, allouche2006assessing}.

Modelos com elevado AUC podem apresentar curvas de resposta biologicamente
incoerentes, forte dependência de viés amostral ou extrapolações inadequadas.
Assim, a avaliação do modelo deve combinar métricas estatísticas e interpretação
ecológica.

\begin{center}
	\fbox{
		\begin{minipage}{0.90\textwidth}
			
			\textbf{Consequências principais.}
			
			Escolher modelos apenas pelo maior AUC pode gerar falsa sensação de desempenho,
			supervalorização de pequenas diferenças estatísticas e interpretações
			ecológicas equivocadas.
			
		\end{minipage}
	}
\end{center}

Boas práticas incluem avaliar simultaneamente AUC, TSS, Índice de Boyce,
matrizes de confusão, curvas de resposta, mapas de adequabilidade e coerência
biogeográfica das predições.

%------------------------------------------------
\section{Erro 4 - Projetar para o futuro sem avaliar extrapolação}

Projeções futuras assumem que a relação estimada entre ocorrência e ambiente
no presente permanece válida sob novas condições climáticas. Entretanto,
cenários futuros podem conter combinações ambientais inexistentes no período de
calibração, resultando em extrapolação ambiental.

Quando a extrapolação não é avaliada, o modelo pode produzir previsões
aparentemente precisas, mas biologicamente frágeis, especialmente em regiões
com climas futuros muito diferentes dos climas atuais
\citep{elith2010functional, owens2013constraints, zurell2020standard}.

\begin{center}
	\fbox{
		\begin{minipage}{0.90\textwidth}
			
			\textbf{Consequências principais.}
			
			Projetar modelos para o futuro sem avaliar extrapolação pode gerar previsões
			biologicamente absurdas, mapas instáveis e interpretações equivocadas de perda
			ou ganho de áreas adequadas.
			
		\end{minipage}
	}
\end{center}

Boas práticas incluem calcular MESS, MOP, mapas de novidade ambiental e
identificar explicitamente áreas onde as projeções estão fora do domínio
ambiental observado durante a calibração.

%------------------------------------------------
\section{Erro 5 - Confundir adequabilidade ambiental com probabilidade de ocorrência}

Modelos de distribuição geralmente produzem índices de adequabilidade
ambiental, isto é, valores que expressam o grau de compatibilidade entre as
condições ambientais de um local e os registros utilizados na calibração do
modelo. Esses valores não devem ser interpretados automaticamente como
probabilidade real de ocorrência da espécie.

A ocorrência efetiva depende também de dispersão, interações bióticas,
história biogeográfica, dinâmica populacional, uso da paisagem e limitações de
amostragem \citep{soberon2007grinnellian, franklin2010mapping,
	peterson2011ecological}.

\begin{sdmwarningbox}
	\textbf{Consequência principal.}
	
	Confundir adequabilidade ambiental com probabilidade real de ocorrência pode
	levar à superestimação da distribuição da espécie e à interpretação incorreta
	de mapas de adequabilidade como mapas absolutos de presença e ausência.
\end{sdmwarningbox}

A interpretação recomendada é tratar os mapas contínuos como gradientes de
suporte ambiental ao modelo, evitando inferências absolutas sem validação
independente ou conhecimento ecológico adicional.

%------------------------------------------------
\section{Erro 6 - Ignorar o viés amostral}

Registros de ocorrência raramente representam amostras aleatórias da
distribuição das espécies. Em muitos casos, estão concentrados próximos a
estradas, rios navegáveis, cidades, centros de pesquisa, unidades de
conservação ou regiões historicamente mais inventariadas.

Quando esse viés não é tratado, o modelo pode aprender o padrão espacial do
esforço de coleta em vez da distribuição ecológica da espécie
\citep{phillips2009sample, kramer2013importance, fourcade2014mapping}.

\begin{center}
	\fbox{
		\begin{minipage}{0.90\textwidth}
			
			\textbf{Consequência principal.}
			
			Ignorar o viés amostral pode fazer com que o modelo reflita onde os
			pesquisadores coletaram mais dados, e não onde a espécie realmente encontra
			condições ambientais adequadas.
			
		\end{minipage}
	}
\end{center}

Boas práticas incluem rarefação espacial, filtros por distância mínima,
avaliação do esforço amostral, uso de background enviesado quando apropriado e
interpretação cautelosa de regiões subamostradas.

\section{Da qualidade metodológica à reprodutibilidade científica}

Evitar erros metodológicos representa apenas uma parte do processo de produção
de conhecimento científico confiável. Mesmo modelos ecologicamente robustos
podem tornar-se difíceis de verificar ou reutilizar quando os procedimentos
analíticos não são adequadamente documentados.

Por esse motivo, a comunidade científica internacional vem enfatizando a
necessidade de incorporar princípios de reprodutibilidade e ciência aberta ao
fluxo de trabalho em SDM. Essa abordagem busca garantir que os resultados
obtidos possam ser reproduzidos, auditados e reutilizados por outros
pesquisadores, fortalecendo a transparência e a credibilidade da pesquisa
ecológica.

%------------------------------------------------
\section{Reprodutibilidade científica em Modelagem de Distribuição de Espécies}
%------------------------------------------------

\subsection{Por que a reprodutibilidade é importante?}

A reprodutibilidade científica constitui um dos pilares fundamentais do método
científico moderno. Um estudo é considerado reprodutível quando outros
pesquisadores conseguem reproduzir os mesmos resultados utilizando os mesmos
dados, métodos analíticos e procedimentos computacionais descritos pelos autores.

Em Modelagem de Distribuição de Espécies (SDM), a reprodutibilidade possui
importância ainda maior devido à grande quantidade de etapas envolvidas no fluxo
de trabalho, incluindo aquisição de dados, filtragem espacial, seleção de
variáveis ambientais, definição da área acessível ($M$), ajuste de modelos,
avaliação estatística, projeções climáticas e análise de incertezas.

Pequenas diferenças nas versões dos softwares, bibliotecas ou algoritmos podem
produzir resultados distintos, dificultando a comparação entre estudos e a
verificação independente das conclusões. Por esse motivo, periódicos
internacionais e iniciativas de ciência aberta recomendam que todos os
procedimentos computacionais sejam documentados e disponibilizados de forma
transparente \citep{wilkinson2016fair, peng2011reproducible,
	sandve2013ten, zurell2020standard}.

%------------------------------------------------
\subsection{Controle do ambiente computacional}
%------------------------------------------------

Uma etapa fundamental para garantir reprodutibilidade consiste em registrar as
versões do software e dos pacotes utilizados durante a análise.

No ambiente R, a função \texttt{sessionInfo()} fornece informações detalhadas
sobre o sistema operacional, versão do R e bibliotecas carregadas durante a
sessão.

\begin{rcode}
	sessionInfo()
\end{rcode}

A saída dessa função deve ser arquivada juntamente com os scripts analíticos,
permitindo que outros pesquisadores reconstruam o ambiente computacional
original.

\begin{sdmpracticebox}
	Inclua sempre a saída de \texttt{sessionInfo()} nos materiais suplementares,
	repositórios GitHub ou documentos de reprodutibilidade do projeto.
\end{sdmpracticebox}

%------------------------------------------------
\subsection{Gerenciamento de dependências}
%------------------------------------------------

Ao longo do desenvolvimento de um projeto, versões de pacotes podem ser
atualizadas, modificadas ou removidas. Isso pode dificultar a reprodução futura
das análises.

O pacote \texttt{renv} permite criar ambientes isolados contendo exatamente as
versões utilizadas durante o desenvolvimento do estudo.

Após instalar os pacotes necessários, recomenda-se criar um snapshot do
ambiente:

\begin{rcode}
	renv::snapshot()
\end{rcode}

Esse comando gera um arquivo denominado \texttt{renv.lock}, que registra todas
as dependências do projeto.

Posteriormente, outro pesquisador poderá reconstruir exatamente o mesmo
ambiente utilizando:

\begin{rcode}
	renv::restore()
\end{rcode}

Essa abordagem reduz significativamente problemas associados à atualização de
bibliotecas e aumenta a confiabilidade da reprodução dos resultados.

%------------------------------------------------
\subsection{Controle de versão com Git}
%------------------------------------------------

O controle de versão permite registrar todas as modificações realizadas em um
projeto ao longo do tempo.

O sistema Git tornou-se o padrão internacional para gerenciamento de código
científico, oferecendo rastreabilidade, documentação histórica e colaboração
entre pesquisadores.

A criação de um repositório local pode ser iniciada com os seguintes comandos:

\begin{shellcode}
	git init
	git add .
	git commit -m "Primeira versão"
\end{shellcode}

Cada nova modificação relevante pode ser registrada por meio de novos
\textit{commits}, criando um histórico completo do desenvolvimento do projeto.

As principais vantagens incluem:

\begin{itemize}
	\item rastreamento de alterações;
	\item recuperação de versões anteriores;
	\item colaboração entre pesquisadores;
	\item documentação automática do fluxo de trabalho;
	\item integração com plataformas de ciência aberta.
\end{itemize}

%------------------------------------------------
\subsection{Disponibilização pública dos dados e scripts}
%------------------------------------------------

A transparência científica exige que dados, scripts e resultados sejam
disponibilizados sempre que possível.

Atualmente existem diversas plataformas especializadas para armazenamento,
compartilhamento e preservação de dados científicos.

\begin{table}[H]
	\centering
	\caption{Principais plataformas utilizadas para compartilhamento de dados e códigos científicos.}
	\label{tab:repositorios_ciencia_aberta}
	
	\begin{tabular}{
			p{0.20\textwidth}
			p{0.70\textwidth}
		}
		
		\toprule
		
		\textbf{Plataforma} &
		\textbf{Finalidade principal}
		
		\\
		\midrule
		
		GitHub &
		Versionamento de código, documentação e desenvolvimento colaborativo.
		
		\\
		
		Zenodo &
		Arquivamento permanente de dados e geração automática de DOI.
		
		\\
		
		Dryad &
		Repositório especializado para conjuntos de dados científicos.
		
		\\
		
		Figshare &
		Compartilhamento de dados, figuras, tabelas e materiais suplementares.
		
		\\
		
		\bottomrule
		
	\end{tabular}
	
\end{table}

Uma prática amplamente recomendada consiste em desenvolver o projeto em um
repositório GitHub e, posteriormente, arquivar uma versão estável no Zenodo,
obtendo um DOI permanente para citação.

%------------------------------------------------
\subsection{Princípios FAIR}
%------------------------------------------------

Em 2016 foi proposto um conjunto de diretrizes conhecido como FAIR Principles,
que se tornou referência internacional para gestão e compartilhamento de dados
científicos \citep{wilkinson2016fair}.

FAIR representa os seguintes princípios:

\begin{table}[H]
	\centering
	\caption{Princípios FAIR para gestão de dados científicos.}
	\label{tab:fair_principles}
	
	\begin{tabular}{
			p{0.15\textwidth}
			p{0.75\textwidth}
		}
		
		\toprule
		
		\textbf{Princípio} &
		\textbf{Descrição}
		
		\\
		\midrule
		
		Findable &
		Dados devem ser facilmente encontrados por humanos e sistemas computacionais.
		
		\\
		
		Accessible &
		Dados devem estar acessíveis por protocolos padronizados.
		
		\\
		
		Interoperable &
		Dados devem utilizar formatos compatíveis e padronizados.
		
		\\
		
		Reusable &
		Dados devem conter documentação suficiente para reutilização futura.
		
		\\
		
		\bottomrule
		
	\end{tabular}
	
\end{table}

A adoção dos princípios FAIR aumenta a transparência, a reutilização dos dados
e o impacto científico dos estudos.

%------------------------------------------------
\subsection{Checklist de reprodutibilidade}
%------------------------------------------------

Antes da publicação de um estudo de SDM, recomenda-se verificar os seguintes
itens:

\begin{itemize}
	\item Os scripts foram disponibilizados?
	\item As versões dos pacotes foram registradas?
	\item O ambiente computacional foi documentado?
	\item Os dados estão disponíveis em repositório público?
	\item O projeto possui DOI?
	\item Os procedimentos metodológicos podem ser reproduzidos integralmente?
	\item Os arquivos possuem metadados adequados?
	\item Os princípios FAIR foram considerados?
\end{itemize}

%------------------------------------------------
\subsection{Mensagem final}
%------------------------------------------------

A reprodutibilidade não deve ser vista como uma etapa adicional do projeto,
mas como parte integrante do processo científico. Em estudos de Modelagem de
Distribuição de Espécies, a documentação adequada dos dados, códigos,
dependências computacionais e decisões metodológicas aumenta a transparência,
fortalece a credibilidade dos resultados e permite que o conhecimento produzido
seja efetivamente reutilizado pela comunidade científica.

\begin{sdmkeybox}
	Um modelo de distribuição somente pode ser considerado plenamente
	científico quando seus resultados podem ser reproduzidos de forma
	independente por outros pesquisadores utilizando os mesmos dados,
	métodos e procedimentos descritos no estudo.
\end{sdmkeybox}

%------------------------------------------------
\section{Síntese dos principais erros e boas práticas}

\begin{table}[H]
	\centering
	\caption{Principais erros em SDM, consequências e boas práticas recomendadas.}
	\label{tab:erros_comuns_sdm}
	
	\scriptsize
	
	\begin{tabular}{
			p{0.24\textwidth}
			p{0.34\textwidth}
			p{0.34\textwidth}
		}
		
		\toprule
		\textbf{Erro comum} &
		\textbf{Consequências} &
		\textbf{Boas práticas} \\
		\midrule
		
		Usar todas as variáveis do WorldClim &
		Multicolinearidade, sobreajuste, projeções instáveis e baixa transferibilidade &
		Correlação, VIF, PCA e seleção ecológica de variáveis \\
		
		Ignorar a área acessível ($M$) &
		Pseudo-ausências irreais, background inadequado e viés de calibração &
		Delimitar $M$ com base em biogeografia, dispersão e barreiras históricas \\
		
		Escolher modelo apenas pelo AUC &
		Falsa sensação de desempenho e interpretações equivocadas &
		Combinar AUC, TSS, Boyce, curvas de resposta e plausibilidade ecológica \\
		
		Projetar para o futuro sem avaliar extrapolação &
		Previsões biologicamente frágeis e extrapolação climática não detectada &
		Usar MESS, MOP e mapas de novidade ambiental \\
		
		Confundir adequabilidade com probabilidade de ocorrência &
		Superestimação da ocupação real e interpretações incorretas &
		Interpretar mapas como gradientes de adequabilidade ambiental \\
		
		Ignorar viés amostral &
		Modelos refletem esforço de coleta, e não distribuição real &
		Rarefação espacial, filtros por distância e background adequado \\
		
		\bottomrule
	\end{tabular}
	
\end{table}

\section{Síntese da unidade}

Ao longo desta unidade foram discutidos alguns dos principais desafios
metodológicos enfrentados em estudos de Modelagem de Distribuição de Espécies.
Foi demonstrado que a qualidade de um modelo depende não apenas do algoritmo
utilizado, mas também da definição adequada da área acessível, da seleção de
variáveis ambientais, do controle do viés amostral, da avaliação da
extrapolação e da interpretação crítica das métricas de desempenho.

Também foram apresentados os fundamentos da reprodutibilidade científica,
controle de ambiente computacional, gerenciamento de dependências, controle de
versão, ciência aberta e princípios FAIR. Esses elementos constituem a base das
boas práticas atualmente recomendadas por periódicos internacionais,
organizações científicas e iniciativas globais de compartilhamento de dados.

Mais do que produzir mapas de adequabilidade ambiental, a modelagem moderna
busca gerar conhecimento transparente, verificável e reutilizável, capaz de
contribuir efetivamente para a compreensão da biodiversidade e para a tomada de
decisões em conservação sob cenários de rápidas mudanças ambientais.

%------------------------------------------------
\section{Mensagem final}

A maioria dos erros em SDM não está relacionada ao algoritmo utilizado, mas às
decisões tomadas antes e depois do ajuste do modelo. A preparação dos dados, a
definição da área acessível, a seleção de variáveis, o controle do viés
amostral, a avaliação da extrapolação e a interpretação ecológica dos
resultados são etapas tão importantes quanto o ajuste estatístico.

\begin{sdmkeybox}
	Um modelo simples, ecologicamente bem fundamentado e cuidadosamente avaliado
	pode ser mais confiável do que um modelo complexo ajustado com muitos
	preditores, área de calibração inadequada, viés amostral não tratado e
	interpretação exclusivamente baseada em métricas estatísticas.
\end{sdmkeybox}

%================================================
